% !TEX root = ../../../main.tex
% 所属：第二部分：多变量分析学 / 第2章：多元函数微分学
% 节：2.2 复合函数的链式法则
% 正文按原文件顺序保存；由 ../chapter2.tex 调用。
\section{复合函数的链式法则}

\subsection*{二元函数复合函数求导法则}

设
\[
f:D_f\subset\R^2\to\R,
\qquad
g:D_g\subset\R^2\to\R^2,
\]
且 $g(D_g)\subset D_f$.记
\[
g(u,v)=\bigl(x(u,v),y(u,v)\bigr).
\]

\begin{theorem}
	若 $g$ 在 $(u_0,v_0)$ 处可微，记
	\[
	x_0=x(u_0,v_0),\qquad y_0=y(u_0,v_0),
	\]
	并且 $f$ 在 $(x_0,y_0)$ 处可微，则复合函数
	\[
	z=f\circ g=f\bigl(x(u,v),y(u,v)\bigr)
	\]
	在 $(u_0,v_0)$ 处可微，且
	\[
	\frac{\partial z}{\partial u}
	=
	\frac{\partial f}{\partial x}\frac{\partial x}{\partial u}
	+
	\frac{\partial f}{\partial y}\frac{\partial y}{\partial u},
	\qquad
	\frac{\partial z}{\partial v}
	=
	\frac{\partial f}{\partial x}\frac{\partial x}{\partial v}
	+
	\frac{\partial f}{\partial y}\frac{\partial y}{\partial v},
	\]
	其中 $\dfrac{\partial f}{\partial x}$、$\dfrac{\partial f}{\partial y}$
	在 $(x_0,y_0)$ 处取值，而
	$\dfrac{\partial x}{\partial u}$、$\dfrac{\partial y}{\partial u}$、
	$\dfrac{\partial x}{\partial v}$、$\dfrac{\partial y}{\partial v}$
	在 $(u_0,v_0)$ 处取值.
\end{theorem}

上述链式法则可写成矩阵形式
\[
\begin{pmatrix}
\dfrac{\partial z}{\partial u}\\[5pt]
\dfrac{\partial z}{\partial v}
\end{pmatrix}_{(u_0,v_0)}
=
\begin{pmatrix}
\dfrac{\partial x}{\partial u}
&\dfrac{\partial y}{\partial u}\\[5pt]
\dfrac{\partial x}{\partial v}
&\dfrac{\partial y}{\partial v}
\end{pmatrix}_{(u_0,v_0)}
\begin{pmatrix}
\dfrac{\partial f}{\partial x}\\[5pt]
\dfrac{\partial f}{\partial y}
\end{pmatrix}_{(x_0,y_0)}.
\]

\begin{proof}
	这里只证明关于 $u$ 的公式.记
	\[
	\begin{aligned}
	\Delta x&=x(u_0+\Delta u,v_0)-x(u_0,v_0),\\
	\Delta y&=y(u_0+\Delta u,v_0)-y(u_0,v_0).
	\end{aligned}
	\]
	于是
	\[
	\begin{aligned}
	\frac{\partial z}{\partial u}(u_0,v_0)
	&=
	\lim_{\Delta u\to0}
	\frac{
	f(x_0+\Delta x,y_0+\Delta y)-f(x_0,y_0)}
	{\Delta u}.
	\end{aligned}
	\]
	由 $f$ 在 $(x_0,y_0)$ 处可微，
	\[
	f(x_0+\Delta x,y_0+\Delta y)-f(x_0,y_0)
	=
	f_x(x_0,y_0)\Delta x
	+
	f_y(x_0,y_0)\Delta y
	+
	o\left(\sqrt{\Delta x^2+\Delta y^2}\right).
	\]
	又因 $g$ 在 $(u_0,v_0)$ 处可微，
	\[
	\lim_{\Delta u\to0}\frac{\Delta x}{\Delta u}
	=
	\frac{\partial x}{\partial u}(u_0,v_0),
	\qquad
	\lim_{\Delta u\to0}\frac{\Delta y}{\Delta u}
	=
	\frac{\partial y}{\partial u}(u_0,v_0),
	\]
	并且
	\[
	\lim_{\Delta u\to0}
	\frac{o\left(\sqrt{\Delta x^2+\Delta y^2}\right)}
	{\Delta u}=0.
	\]
	故
	\[
	\frac{\partial z}{\partial u}
	=
	\frac{\partial f}{\partial x}\frac{\partial x}{\partial u}
	+
	\frac{\partial f}{\partial y}\frac{\partial y}{\partial u}.
	\]
	关于 $v$ 的公式同理.
\end{proof}

一般地，设
\[
f:D_f\subset\R^m\to\R,
\qquad
g:D_g\subset\R^n\to\R^m,
\]
并记
\[
g(x_1,\cdots,x_n)
=
\bigl(y_1(x_1,\cdots,x_n),\cdots,y_m(x_1,\cdots,x_n)\bigr),
\qquad
z=f\circ g.
\]
若 $f$ 与 $g$ 在相应点处可微，则对于 $i=1,\cdots,n$，
\[
\frac{\partial z}{\partial x_i}
=
\frac{\partial f}{\partial y_1}\frac{\partial y_1}{\partial x_i}
+
\frac{\partial f}{\partial y_2}\frac{\partial y_2}{\partial x_i}
+\cdots+
\frac{\partial f}{\partial y_m}\frac{\partial y_m}{\partial x_i}.
\]
写成矩阵形式即
\[
\begin{pmatrix}
\dfrac{\partial z}{\partial x_1}
&\cdots&
\dfrac{\partial z}{\partial x_n}
\end{pmatrix}
=
\begin{pmatrix}
\dfrac{\partial f}{\partial y_1}
&\cdots&
\dfrac{\partial f}{\partial y_m}
\end{pmatrix}
\frac{D(y_1,\cdots,y_m)}
{D(x_1,\cdots,x_n)}.
\]

\begin{theorem}[向量值复合函数的链式法则]
	设
	\[
	\vect f:D_f\subset\R^k\to\R^m,
	\qquad
	\vect g:D_g\subset\R^n\to\R^k
	\]
	分别为多元向量值函数，且分别在 $D_f$ 与 $D_g$ 上具有连续导数.
	若 $\vect g(D_g)\subset D_f$，并记 $\vect u=\vect g(\vect x)$，
	则复合向量值函数 $\vect f\circ\vect g$ 在 $D_g$ 上也具有连续导数，
	并且
	\[
	(\vect f\circ\vect g)'(\vect x)
	=
	\vect f'(\vect u)\,\vect g'(\vect x)
	=
	\vect f'\bigl(\vect g(\vect x)\bigr)\,\vect g'(\vect x).
	\]
	其中 $\vect f'(\vect u)$、$\vect g'(\vect x)$ 和
	$(\vect f\circ\vect g)'(\vect x)$ 分别表示相应函数的 Jacobi 矩阵.
\end{theorem}

\subsection*{例子}

\begin{example}
	设
	\[
	z=\arctan(xy),\qquad y=e^x.
	\]
	求 $\left.\dfrac{\dd z}{\dd x}\right|_{x=0}$.

	由复合函数的链式法则，
	\[
	\frac{\dd z}{\dd x}
	=
	\frac{\partial z}{\partial x}
	+
	\frac{\partial z}{\partial y}\frac{\dd y}{\dd x}
	=
	\frac{y}{1+x^2y^2}
	+
	\frac{x}{1+x^2y^2}e^x.
	\]
	当 $x=0$ 时，$y=1$，故
	\[
	\left.\frac{\dd z}{\dd x}\right|_{x=0}=1.
	\]
\end{example}

\begin{example}
	设
	\[
	z=\frac{x^2}{y},\qquad x=2uv,\qquad y=2u+v.
	\]
	则
	\[
	\begin{aligned}
	\frac{\partial z}{\partial u}
	&=
	\frac{\partial z}{\partial x}\frac{\partial x}{\partial u}
	+
	\frac{\partial z}{\partial y}\frac{\partial y}{\partial u}\\
	&=
	\frac{2x}{y}\cdot 2v-\frac{x^2}{y^2}\cdot 2\\
	&=
	\frac{4uv}{2u+v}\cdot 2v
	-
	\frac{8u^2v^2}{(2u+v)^2}.
	\end{aligned}
	\]
\end{example}

\begin{example}
	设
	\[
	w=f(x^2+y^2,xyz).
	\]
	令
	\[
	\xi=x^2+y^2,
	\qquad
	\eta=xyz,
	\qquad
	w=f(\xi,\eta).
	\]
	则
	\[
	\frac{\partial w}{\partial x}
	=
	\frac{\partial f}{\partial\xi}\frac{\partial\xi}{\partial x}
	+
	\frac{\partial f}{\partial\eta}\frac{\partial\eta}{\partial x}
	=
	2x\frac{\partial f}{\partial\xi}
	+yz\frac{\partial f}{\partial\eta}.
	\]
	这里 $\dfrac{\partial f}{\partial\xi}$ 称为 $f$ 对第一变元求偏导，
	$\dfrac{\partial f}{\partial\eta}$ 称为 $f$ 对第二变元求偏导.
	在所需的二阶偏导数均存在时，进一步有
	\[
	\begin{aligned}
	\frac{\partial^2w}{\partial z\partial x}
	&=
	\frac{\partial}{\partial z}
	\left(
	2x\frac{\partial f}{\partial\xi}
	+yz\frac{\partial f}{\partial\eta}
	\right)\\
	&=
	2x\left(
	\frac{\partial^2f}{\partial\xi^2}\frac{\partial\xi}{\partial z}
	+
	\frac{\partial^2f}{\partial\eta\partial\xi}
	\frac{\partial\eta}{\partial z}
	\right)
	+y\frac{\partial f}{\partial\eta}\\
	&\qquad
	+yz\left(
	\frac{\partial^2f}{\partial\xi\partial\eta}
	\frac{\partial\xi}{\partial z}
	+
	\frac{\partial^2f}{\partial\eta^2}
	\frac{\partial\eta}{\partial z}
	\right)\\
	&=
	2x^2y\frac{\partial^2f}{\partial\eta\partial\xi}
	+y\frac{\partial f}{\partial\eta}
	+xy^2z\frac{\partial^2f}{\partial\eta^2}.
	\end{aligned}
	\]
\end{example}

\subsection*{一阶微分形式的不变性}

回忆一元复合函数的情形.设
\[
y=f(x),\qquad x=\varphi(t).
\]
若把 $x$ 看作中间变量，对 $t$ 求微分，则
\[
\dd y=f'(x)\varphi'(t)\dd t=f'(x)\dd x;
\]
若直接把 $x$ 看作自变量，同样有
\[
\dd y=f'(x)\dd x.
\]

类似地，设
\[
z=f(x,y),
\qquad
x=x(u,v),
\qquad
y=y(u,v).
\]
当 $x,y$ 被看作自变量时，
\[
\dd z=\frac{\partial f}{\partial x}\dd x
+\frac{\partial f}{\partial y}\dd y.
\]
当 $x,y$ 被看作中间变量时，
\[
\begin{aligned}
\dd z
&=
\frac{\partial z}{\partial u}\dd u
+
\frac{\partial z}{\partial v}\dd v\\
&=
\left(
\frac{\partial f}{\partial x}\frac{\partial x}{\partial u}
+
\frac{\partial f}{\partial y}\frac{\partial y}{\partial u}
\right)\dd u
+
\left(
\frac{\partial f}{\partial x}\frac{\partial x}{\partial v}
+
\frac{\partial f}{\partial y}\frac{\partial y}{\partial v}
\right)\dd v\\
&=
\frac{\partial f}{\partial x}
\left(
\frac{\partial x}{\partial u}\dd u
+
\frac{\partial x}{\partial v}\dd v
\right)
+
\frac{\partial f}{\partial y}
\left(
\frac{\partial y}{\partial u}\dd u
+
\frac{\partial y}{\partial v}\dd v
\right)\\
&=
\frac{\partial f}{\partial x}\dd x
+
\frac{\partial f}{\partial y}\dd y.
\end{aligned}
\]
因此，无论 $x,y$ 是自变量还是中间变量，一阶全微分都具有相同的形式，
这称为\textbf{一阶微分形式的不变性}.一般而言，二阶及更高阶微分形式不具有这种不变性.

